% preamble.tex —— 《点燃物理》全局样式
\usepackage{amsmath,amssymb,bm}
\usepackage{graphicx}
\usepackage{booktabs}
\usepackage{enumitem}
\usepackage{tikz}
\usetikzlibrary{arrows.meta,positioning,calc,decorations.pathmorphing,patterns}
\usepackage[most]{tcolorbox}
\usepackage[a4paper,margin=2.4cm,top=2.6cm,bottom=2.8cm]{geometry}
\usepackage{hyperref}
\hypersetup{colorlinks=true,linkcolor=black,urlcolor=blue!50!black,citecolor=black}

% ---------- 三层阅读标记 ----------
% 主线层：只要求初中数学，保证故事完整
\newtcolorbox{mainline}{%
  colback=green!4,colframe=green!50!black,breakable,
  title=【主线层】,fonttitle=\bfseries,
  left=6pt,right=6pt,top=4pt,bottom=4pt}
% 数学桥层：引入向量、函数、极限、导数、积分和简单微分方程
\newtcolorbox{mathbridge}{%
  colback=blue!4,colframe=blue!55!black,breakable,
  title=【数学桥层】,fonttitle=\bfseries,
  left=6pt,right=6pt,top=4pt,bottom=4pt}
% 挑战层：线性代数、偏微分方程、复数等，非理解后文的硬门槛
\newtcolorbox{challenge}{%
  colback=red!4,colframe=red!60!black,breakable,
  title=【挑战层】,fonttitle=\bfseries,
  left=6pt,right=6pt,top=4pt,bottom=4pt}

% ---------- 固定结构盒子 ----------
% 动手实验
\newtcolorbox{experiment}[1][]{%
  colback=orange!6,colframe=orange!75!black,breakable,
  title={动手实验\ifx\\#1\\\else：#1\fi},fonttitle=\bfseries,
  left=6pt,right=6pt,top=4pt,bottom=4pt}
% 脑洞实验 / 思想实验
\newtcolorbox{thought}[1][]{%
  colback=violet!6,colframe=violet!70!black,breakable,
  title={脑洞实验\ifx\\#1\\\else：#1\fi},fonttitle=\bfseries,
  left=6pt,right=6pt,top=4pt,bottom=4pt}
% 一句话模型
\newtcolorbox{oneline}{%
  colback=yellow!10,colframe=yellow!60!black,breakable,
  title=一句话模型,fonttitle=\bfseries,
  left=6pt,right=6pt,top=4pt,bottom=4pt}
% 写作提纲占位（仅骨架章节使用，成书时删除）
\newtcolorbox{draftnote}{%
  colback=gray!6,colframe=gray!60,breakable,
  title={本章写作提纲（占位，成书时删除）},fonttitle=\bfseries,
  left=6pt,right=6pt,top=4pt,bottom=4pt}

% ---------- 每章十步固定结构的小节标题 ----------
% 用法：\step{反常识开场}
\newcommand{\step}[1]{\subsection*{#1}}
% 挑战题
\newlist{puzzles}{enumerate}{1}
\setlist[puzzles]{label=\textbf{挑战 \arabic*.},leftmargin=3em}

% ---------- 册（三册） ----------
% 用法：\volume{第一册}{运动、能量与场}
\newcommand{\volume}[2]{%
  \cleardoublepage
  \thispagestyle{empty}
  \vspace*{0.28\textheight}
  \begin{center}
    {\Huge\bfseries #1\par}
    \vspace{1.2em}
    {\LARGE #2\par}
  \end{center}
  \vfill
  \phantomsection
  \addcontentsline{toc}{part}{#1\quad #2}
  \cleardoublepage}

% ---------- 篇（十一篇） ----------
% 直接用 ctexbook 的 \part，去掉自动编号重设即可（book 类默认不重置章编号）
